% Ad Familiares 1.8 (ad-familiares-01-08) --- English translation
% Translated by Claude (Anthropic) from the Latin (Perseus).
% Style guide: STYLE.md.

\ciceroLetterOpener{M. CICERO S. D. P. LENTVLO PROCOS.}

\ciceroSection{1} On all matters that touch you --- what has been done, what has been settled, what Pompey has taken in hand --- you will best learn from Marcus Plaetorius, who has not only had a part in these matters but has been at their head, and has omitted no service that the most affectionate, the most prudent, the most attentive of men could render you. From the same man you will learn the whole condition of public affairs. What kind of condition that is, is not easy to write. Certainly they are now in the power of our friends, and so much so that no change in this generation seems possible.

\ciceroSection{2} For my own part, as duty requires, as you yourself instructed me, and as obligation and interest compel me, I am attaching myself to the policy of one whom you yourself thought it useful to attach to your own policy. But you do not need to be told how hard it is to lay down a feeling about politics, especially one well grounded and confirmed. Yet I am bringing myself round to the will of the man from whom I cannot honourably differ; and I am not doing so, as some perhaps may think, in pretence: so strong with me is the bent of my mind --- and, by Hercules, my affection for Pompey --- that whatever serves him and whatever he wills now seem to me both right and true; nor, in my judgement, would even his opponents go astray, if, since they cannot match him, they would stop fighting.

\ciceroSection{3} Even this consoles me: that I am the man to whom all would most readily yield the right either to defend whatever Pompey wishes, or to be silent, or even --- which I most prefer --- to come back to my own studies of letters. And come back to them I shall, surely, if I can do so by his friendship. For the prospect that lay before me, when I had run through the highest offices and the greatest labours --- standing in the giving of opinions, freedom in the conduct of the state --- has been swept away entirely, and not for me more than for everyone else: for now one must either assent without dignity along with the few, or dissent in vain.

\ciceroSection{4} I write this to you above all for this reason, that you may now think over your own policy too. The whole order of the Senate, of the courts, of the entire commonwealth, has been changed: peace is what we must wish for, which the men who hold the power seem ready to provide --- if certain people will be able to bear their power more patiently. As for that consular dignity of the brave and constant senator, there is no thinking of it: that has been lost, by the fault of those who alienated from the Senate both an order most closely bound to it and a man of the highest distinction.

\ciceroSection{5} But to come back to what is more closely bound up with your own affairs: I have learned that Pompey is your warm friend; and with him as consul, so far as I can see, you will obtain everything you wish; in which matters he will have me fastened to him, nor will any matter touching you be neglected by me; for I shall not be afraid of being a burden to one who will be glad of it on his own account, when he sees that I am a grateful man.

\ciceroSection{6} I should like you to be assured that there is no matter, however small, touching you, that is not dearer to me than all my own affairs; and although in this conviction I can satisfy myself by my zeal, in the matter itself I cannot, for the reason that I cannot reach by repaying --- nor even by holding in mind --- any portion of your services to me.

\ciceroSection{7} The rumour was that you have managed your affairs very well; your dispatches were being awaited, of which we had already spoken with Pompey. When they have arrived, our zeal will show itself in addressing the magistrates and the senators; and in the rest of what touches you, even when we shall have exerted ourselves more than we can, we shall still do less than we ought.
